\section{Modifications to chapter 13}

The following modifications shall be applied to chapter 13 (\emph{Templates}
[\texttt{temp}]):

\begin{modification}[Modification to 13.7.1/3]
A \emph{separately instantiated construct} of a templated function \texttt{F} is a

\begin{itemize}[label=\textemdash]

\item
default argument,

\item
\emph{noexcept-specifier} , or

\item
\stdremove{\emph{function-contract-specifier}}

\end{itemize}

of \texttt{F}. For purposes of name lookup and instantiation, separately
instantiated constructs, \emph{type-constraints}, and
\emph{requires-clauses} (13.1) are considered definitions; each separately
instantiated construct, type-constraint, or \emph{requires-clause} is a
separate definition which is unrelated to the templated function definition or
to any other separately instantiated constructs, \emph{type-constraints}, or
\emph{requires-clauses}. For the purpose of instantiation, the substatements of
a constexpr if statement (8.5.2) are considered definitions. For the purpose of
name lookup and instantiation, the \emph{compound-statement} of an
\emph{expansion-statement} is considered a template definition.
\end{modification}

\begin{modification}[Modification of 13.8.3.3/3]
An \emph{id-expression} is type-dependent if it is a \emph{template-id} that is
not a \emph{concept-id} and is dependent; or if its
terminal name is

\begin{itemize}[label=\textemdash]

\item \ldots

\item
\stdremove{associated by name lookup with a result binding (9.4.2) of a function whose
return type is dependent,}
\item \ldots

\end{itemize}

\end{modification}

\begin{modification}[Modification of 13.9.4/12]
Whether an explicit specialization of a function or variable template is inline,
constexpr, constinit, or consteval is determined by the explicit specialization
and is independent of those properties of the template. Similarly, attributes
\stdremove{and \emph{function-contract-specifiers}} appearing in the declaration
of a template have no effect on an explicit specialization of that template.
\end{modification}

\begin{modification}[Modification of 13.10.3.1/7]
\ldots\\
NOTE 4 The equivalent substitution in exception specifications (14.5)
\stdremove{and function contract assertions (9.4.1)} is done
only when the \emph{noexcept-specifier} \stdremove{or
function-contract-specifier , respectively,} is instantiated, at which point a program
is ill-formed if the substitution results in an invalid type or expression.
\end{modification}


